\documentclass[11pt, a4paper]{article}

\usepackage[T1]{fontenc}
\usepackage[utf8]{inputenc}
\usepackage{tgpagella}
\usepackage{microtype}
\usepackage[spanish, es-noshorthands]{babel}
\usepackage{amsmath, amssymb}
\usepackage[most]{tcolorbox}
\usepackage[margin=2.5cm]{geometry}
\usepackage{fancyhdr}

\pagestyle{fancy}
\fancyhf{}
\setlength{\headheight}{16pt}
\lhead{\textbf{IMT-342}}
\chead{Handoff de Contenido: Dinámica}
\rhead{Preparación para W09-S02}

\definecolor{ucbblue}{RGB}{0, 51, 102}

\begin{document}

\begin{center}
    \Large\textbf{\textcolor{ucbblue}{Handoff Técnico para W09-S02 (Dinámica)}}\\[0.2cm]
\end{center}

\begin{tcolorbox}[colback=red!5, colframe=red!70!black, title=\textbf{Objetivo Principal de W09-S02[cite: 13]}]
    Introducir la \textbf{estructura dinámica mínima y rigurosa} necesaria para que el estudiante defienda el EC2 y conecte el Hito 2. NO se debe intentar realizar una derivación simbólica de los 6 grados de libertad del robot IRB120, a menos que se suministre explícitamente información geométrica inercial del fabricante.
\end{tcolorbox}

\section*{1. Ecuación Central de Diseño[cite: 13]}
La próxima sesión debe explicar e interpretar la formulación:
\begin{equation*}
    \tau = M(\mathbf{q})\ddot{\mathbf{q}} + C(\mathbf{q}, \dot{\mathbf{q}})\dot{\mathbf{q}} + g(\mathbf{q})
\end{equation*}

\section*{2. Requerimientos de Comprensión del Estudiante[cite: 13]}
El estudiante debe comprender y nombrar el aporte de cada término:
\begin{itemize}
    \item $M(\mathbf{q})\ddot{\mathbf{q}}$: Contribución inercial.
    \item $C(\mathbf{q}, \dot{\mathbf{q}})\dot{\mathbf{q}}$: Efectos de Coriolis y centrífugos dependientes de la velocidad.
    \item $g(\mathbf{q})$: Contribución de compensación de gravedad.
    \item $\tau$: Vector de torque final requerido de los actuadores.
\end{itemize}

\section*{3. Alcance de la Derivación[cite: 13]}
Desarrolle una derivación analítica \textbf{manejable}, por ejemplo, un péndulo de 1R o un brazo planar de 2R. El enfoque es procedimental.

\section*{4. Conexión Obligatoria con el Hito 2 (Trayectoria)[cite: 13]}
Conecte explícitamente la dinámica con W09-S01: Utilice $q(t), \dot{q}(t), \ddot{q}(t)$ generados por el perfil LSPB para inyectarlos en la ecuación de $\tau$.

\section*{5. Salida Computacional Esperada[cite: 13]}
En Python, el estudiante deberá evaluar $\tau(t)$ a lo largo de toda la trayectoria e interpretar los picos de esfuerzo mecánico relativos a las tres fases del LSPB (aceleración, crucero, desaceleración).

\end{document}
